% ----------------------------------------------------------------------------
%                               METODOLOGIA
% ----------------------------------------------------------------------------

\chapter{METODOLOGIA}
\label{chap:metodologia}

O arquivo de partida é denominado {\ttfamily principal.tex}. Este \emph{template} está organizado em três diretórios principais: a pasta {\ttfamily /configuracoes}, que contém configurações relacionadas à classe abn\TeX{}2, configurações de saída do PDF e a lista de pacotes empregados; o diretório {\ttfamily /dados}, onde é recomendado organizar Algoritmos, Figuras, Quadros e Tabelas utilizados no trabalho; e a pasta {\ttfamily /estrutura}, que possui os arquivos onde o TCC deve ser redigido.

Aqui conterão os métodos e procedimentos adotados no desenvolvimento do trabalho. Esta é uma das sessões mais importantes pois demonstra o poder científico que foi utilizado para a pesquisa. Sem uma boa metodologia a pesquisa pode perder a validade. O pesquisador deve utilizar métodos ou técnicas aceitas pela comunidade científica na busca de provar suas hipóteses.

A metodologia escolhida deve ser aquela que mais se adéqua ao seu objeto de estudo e à abordagem aplicada. Há dois métodos principais: 1) quantitativo, que é o uso de instrumental estatístico, de dados numéricos; e 2) qualitativo, que se caracteriza pela qualificação dos dados coletados, durante a análise do problema.

\section{Uma seção}
\lipsum[1]


\subsection{Uma subseção}
\lipsum[1]